\scanpage{324}\mspage{48}
% archive sheet number: 323

\diagram{\begin{tikzcd}[ampersand replacement=\&,row sep=1ex,column sep=1em]
 \& \overset{\unsure{\mathcal{F}|F}_{W_0},\,\mathrm{Fibf}_{W_0}}{\overbrace{A_{W_0}}}
  \arrow[dd,Rightarrow] \arrow[ddddrr,Rightarrow,postaction={draw,decorate,decoration={zigzag,segment length=4pt,amplitude=2pt}},"{\deleted{\diagup}}" description] \& \& \&\\
 \overset{\mathcal{F}_{W},\,\mathrm{Fibf}_{W}}{\overbrace{A}}
  \arrow[ur,Rightarrow] \arrow[dd,Rightarrow] \& \& \& \&\\
 \& B_{W_0}\arrow[d,Rightarrow] \& \& \&\\
 B\arrow[ur,Rightarrow]\arrow[dd,Rightarrow]
 \& C_{W_0}\arrow[uuu,Rightarrow,dashed,bend right=37,"{\text{si }W_0(7)}"']
 \arrow[drr,Rightarrow,"\bullet"] \& \& \&\\
 \& \& \& \overset{\mathrm{Parf}_{W_0}}{\overbrace{D_{W_0}}}
 \arrow[r,Rightarrow,"\bullet"]
 \& \substack{F:\\(\mathrm{Cat},W,\mathrm{Fib}_{W})\\\text{une catégorie}\\\text{à fibrations},}\\
 C\arrow[uur,Rightarrow]\arrow[rr,Rightarrow,"\bullet"] \& \&
 \overset{\mathrm{Parf}_{W}}{\overbrace{D}}\arrow[ur,Rightarrow] \& \&
\end{tikzcd}}

\diagram{\begin{tikzcd}[ampersand replacement=\&,row sep=.35em,column sep=.8em]
 \& A\arrow[r,Rightarrow]\arrow[d,Rightarrow] \& A_{W_0}\arrow[d,Rightarrow] \&\\
 \mathcal{F}_{W},\,\mathrm{Fibf}_{W}\;\smash{\left\{\vphantom{\rule{0pt}{5ex}}\right.}
 \& B\arrow[r,Rightarrow]\arrow[d,Rightarrow] \& \unsure{B|D}_{W_0}\arrow[d,Rightarrow] \&
 \smash{\left.\vphantom{\rule{0pt}{5ex}}\right\}}\;\substack{\cdot\,\cdot\,\cdot\\[-.25ex]\cdot\,\cdot}\\
 \& C\arrow[r,Rightarrow]\arrow[d,Rightarrow] \& C_{W_0}\arrow[d,Rightarrow] \&\\
 \& \underset{\mathrm{Parf}_{W}}{\underbrace{D}}\arrow[r,Rightarrow]
 \& \underset{\mathrm{Parf}_{W_0}}{\underbrace{D_{W_0}}}\arrow[r,Rightarrow] \& F
\end{tikzcd}}
% slanted
\margin{version revue du diagramme}

% slanted
\margin{\deleted{3ème mouture du diagramme}}
\diagram{\begin{tikzcd}[ampersand replacement=\&,row sep=.45em,column sep=.65em,execute at end picture={\draw[overlay,thin] ([xshift=1mm]current bounding box.south west) -- ([xshift=8mm]current bounding box.north west) ([xshift=9mm]current bounding box.south west) -- ([xshift=17mm]current bounding box.north west) ([xshift=17mm]current bounding box.south west) -- ([xshift=23mm]current bounding box.north west);}]
 \deleted{A}\arrow[d,dash]\arrow[r,dash] \& \deleted{A_{W_0}}\arrow[d,dash] \& \&\\
 \deleted{B}\arrow[d,dash]\arrow[r,dash] \& \deleted{B_{W_0}}\arrow[d,dash] \& \&\\
 \deleted{C}\arrow[r,dash]\arrow[dr,dash] \& \deleted{C_{W_0}}\arrow[dr,dash] \& \&\\
 \& \deleted{D}\arrow[r,dash] \& \deleted{D_{W_0}} \& \deleted{F}\arrow[l,"{\unsure{\diagup}}" description]
\end{tikzcd}}

diagramme que je vais commenter. Toutes les implications \inserted{en traits pleins} du diagramme ont \inserted{été déjà \unsure{obtenues|établies} ou sont} tautologiques, à l'exception des trois $\overset{\bullet}{\Longrightarrow}$ sur lesquelles je vais revenir. L'implication en pointillés $C_{W_0}\dashrightarrow A_{W_0}$ est une implication conditionnelle, soumise à la condition marquée sur le diagramme comme $W_0(7)$, ce qui signifie: $W_0$ satisfait à $\mathrm{Loc}(7)$. Cette implication, i.e. l'équivalence de $A_{W_0},B_{W_0},C_{W_0}$ sous la condition $W_0(7)$, a été vue dans 4.4.6. J'ai marqué, au dessus des deux groupes $(A,B,C)$ et $(A_{W_0},B_{W_0},C_{W_0})$, ainsi qu'au dessus de $D$ et de $D'$, qui correspondent chacun à un tel groupe de conditions $A\deleted{\unsure{0}}(\mathcal{F}),B(\mathcal{F}),C(\mathcal{F})$ toutes équivalentes, avec $\mathcal{F}=\deleted{\unsure{\mathcal{F}}}\mathrm{Parf}_{W}$ ou $\mathcal{F}=\deleted{\unsure{\mathcal{F}}}\mathrm{Parf}_{W_0}$\unsure{),}, les
